\documentclass[a4paper,11pt]{article}
\usepackage[utf8]{inputenc}
\usepackage[T1]{fontenc}
\usepackage[french]{babel}
\usepackage{amsmath, amssymb, amsthm}
\usepackage{graphicx}
\usepackage{booktabs}
\usepackage{geometry}
\usepackage{hyperref}
\usepackage{caption}
\usepackage{subcaption}
\usepackage{float}
\usepackage{array}
\usepackage{xcolor}
\usepackage{enumitem}
\usepackage{titlesec}

\geometry{margin=2.2cm}
\hypersetup{colorlinks=true, linkcolor=black, urlcolor=blue}

% ── Mise en forme des titres ──────────────────────────────────
\titleformat{\section}{\large\bfseries}{}{0em}{}[\titlerule]
\titleformat{\subsection}{\normalsize\bfseries}{}{0em}{}

\title{
  \vspace{-1cm}
  {\small Université — Mathématiques pour le Machine Learning}\\[0.4cm]
  \rule{\linewidth}{0.5pt}\\[0.4cm]
  {\LARGE\bfseries Classification d'images par réseaux de neurones}\\[0.2cm]
  {\large De la reconnaissance de chiffres manuscrits\\
          à la détection de cancers du sein}\\[0.3cm]
  \rule{\linewidth}{0.5pt}
}
\author{
  Nazim Mekideche \quad Armand Lauener \\
  Mohamed Ouadihi \quad Macéo Pierson
}
\date{Juin 2026}

\begin{document}

\maketitle
\tableofcontents
\newpage

% ============================================================
\section{Introduction}
% ============================================================

Ce rapport présente l'implémentation et l'analyse mathématique de modèles
de classification d'images progressivement plus complexes, dans le cadre du cours
\textit{Mathématiques pour le Machine Learning}.

Trois jeux de données sont étudiés dans l'ordre croissant de difficulté :
\begin{enumerate}[leftmargin=*]
  \item \textbf{MNIST} — chiffres manuscrits $28\times28$ pixels, 10 classes, 60\,000 images.
  \item \textbf{CIFAR-10} — objets naturels $32\times32$ pixels RGB, 10 classes, 50\,000 images.
  \item \textbf{CBIS-DDSM} — mammographies médicales $224\times224$ pixels, classification
        binaire bénin/malin, 1\,318 images d'entraînement.
\end{enumerate}

Chaque partie introduit de nouvelles contraintes mathématiques :
MNIST illustre les fondements de la classification supervisée ;
CIFAR-10 motive les réseaux convolutifs ;
CBIS-DDSM impose des métriques médicales asymétriques et une gestion du déséquilibre de classes.

% ============================================================
\section{Fondements mathématiques}
% ============================================================

\subsection{Représentation des données et notations}

Une image est représentée par un vecteur $\vec{x} \in \mathbb{R}^d$
(après aplatissement si nécessaire) ou par un tenseur $\mathbf{X} \in \mathbb{R}^{C \times H \times W}$.
Le label associé est encodé en vecteur one-hot $\vec{y} \in \{0,1\}^K$ avec $K$ classes.

\subsection{Fonction softmax et prédiction}

Le vecteur de scores bruts $\vec{z} = A\vec{x} + \vec{b}$ est transformé en distribution
de probabilités par la softmax :
\[
  P_k(\vec{x}) = \frac{e^{z_k}}{\displaystyle\sum_{j=1}^{K} e^{z_j}}, \qquad k \in \{1,\ldots,K\}
\]
La prédiction est $\hat{y} = \arg\max_k P_k(\vec{x})$.

\subsection{Fonction de coût — entropie croisée}

La perte log-loss (entropie croisée catégorielle) pénalise l'écart entre
la distribution prédite et le label réel :
\[
  \mathcal{L}(\vec{y}, P) = -\frac{1}{n}\sum_{i=1}^{n}\sum_{k=1}^{K}
  y_{ik}\ln P_k(\vec{x}_i)
\]
Pour une classification binaire, la BCE (\textit{Binary Cross-Entropy}) utilise
une seule sortie $z$ et la sigmoïde $\sigma(z) = 1/(1+e^{-z})$ :
\[
  \mathcal{L}_{BCE} = -\frac{1}{n}\sum_{i=1}^{n}
  \Bigl[w^+\,y_i\ln\sigma(z_i) + (1-y_i)\ln(1-\sigma(z_i))\Bigr]
\]
où $w^+$ est le poids de la classe positive (maligne en médecine).

\subsection{Rétropropagation et descente de gradient}

Les paramètres $\theta$ sont mis à jour par descente de gradient stochastique (SGD)
avec mini-batches de taille $B$ :
\[
  \theta \leftarrow \theta - \eta\,\nabla_\theta \mathcal{L}(\mathcal{B})
\]
La rétropropagation calcule $\nabla_\theta \mathcal{L}$ par la règle de dérivation composée.
Pour un réseau à $L$ couches, les gradients de la couche $l$ sont :
\[
  \delta^{(l)} = \left(W^{(l+1)}\right)^\top \delta^{(l+1)} \odot \sigma'\!\left(z^{(l)}\right)
\]

\subsection{Fonctions d'activation}

\begin{table}[H]
\centering
\caption{Fonctions d'activation testées dans ce projet.}
\begin{tabular}{lllll}
\toprule
\textbf{Nom} & \textbf{Formule} & \textbf{Dérivée} & \textbf{Image} & \textbf{Remarque} \\
\midrule
Sigmoïde   & $\sigma(z)=\frac{1}{1+e^{-z}}$ & $\sigma(1-\sigma)$ & $(0,1)$ & Gradient qui disparaît en profondeur \\
Tanh       & $\tanh(z)$      & $1-\tanh^2(z)$  & $(-1,1)$ & Centré, mieux que sigmoïde \\
ReLU       & $\max(0,z)$     & $\mathbf{1}_{z>0}$ & $[0,+\infty)$ & Gradient constant, efficace en profondeur \\
Heaviside  & $\mathbf{1}_{z\geq0}$ & $0$ p.p. & $\{0,1\}$ & Non différentiable, gradient nul \\
\bottomrule
\end{tabular}
\end{table}

\noindent\textbf{Problème du gradient qui disparaît :} avec sigmoid/tanh, le gradient
$|\sigma'(z)| \leq 0{,}25$ — après $L$ couches, le gradient est multiplié par
$0{,}25^L \to 0$. ReLU évite ce problème car $\sigma'(z) = 1$ pour $z > 0$.

% ============================================================
\section{Partie 1 — MNIST : classification de chiffres manuscrits}
% ============================================================

\subsection{Description du jeu de données}

MNIST contient $60\,000$ images d'entraînement et $10\,000$ de test,
réparties équitablement sur 10 classes (chiffres 0--9).
Chaque image est un vecteur $\vec{x} \in \mathbb{R}^{784}$ de niveaux de gris normalisés dans $[0,1]$.

\subsection{Modèle linéaire}

Le modèle le plus simple calcule des scores linéaires :
\[
  \vec{z} = A\vec{x} + \vec{b}, \qquad A \in \mathbb{R}^{10 \times 784},\; \vec{b} \in \mathbb{R}^{10}
\]
\textbf{Nombre de paramètres :} $784 \times 10 + 10 = \mathbf{7\,850}$

Ce modèle apprend un ``gabarit'' par classe. Il n'exploite aucune interaction entre pixels,
ce qui limite ses performances sur des formes spatiales complexes.

\subsection{MLP à une couche cachée (H=1)}

On ajoute une couche cachée de $p_1 = 64$ neurones avec activation sigmoid :
\[
  \vec{z}^{(1)} = \sigma\!\left(A^{(1)}\vec{x} + \vec{b}^{(1)}\right) \in \mathbb{R}^{64}, \qquad
  \vec{z}^{(2)} = A^{(2)}\vec{z}^{(1)} + \vec{b}^{(2)} \in \mathbb{R}^{10}
\]
\textbf{Nombre de paramètres :}
$784\times64 + 64 + 64\times10 + 10 = \mathbf{50\,890}$

\textbf{Choix de 64 neurones :} compromis entre capacité (trop peu → sous-apprentissage)
et temps de calcul. La couche cachée permet au réseau de construire des représentations
intermédiaires non linéaires : combinaisons de pixels évoquant des parties de chiffres.

\subsection{MLP à deux couches cachées (H=2)}

On ajoute une seconde couche de $p_2 = 32$ neurones :
\[
  \vec{z}^{(2)} = \sigma\!\left(A^{(2)}\vec{z}^{(1)} + \vec{b}^{(2)}\right) \in \mathbb{R}^{32}
\]
\textbf{Nombre de paramètres :}
$784\times64+64 + 64\times32+32 + 32\times10+10 = \mathbf{52\,650}$

L'architecture en entonnoir $784 \to 64 \to 32 \to 10$ force une compression progressive :
H=1 détecte des caractéristiques simples (bords, courbes), H=2 les combine en motifs complexes.

\subsection{Comparaison des activations × architectures}

Au-delà du minimum requis, nous avons conduit une grille d'expériences
croisant 4 activations (sigmoid, tanh, relu, heaviside) et 5 architectures MLP.
Les principaux enseignements :

\begin{itemize}[leftmargin=*]
  \item \textbf{Heaviside} : gradient nul presque partout — le réseau n'apprend pas
        au-delà de H=1 (gradient bloqué dès la deuxième couche cachée).
  \item \textbf{Sigmoid vs ReLU} : sur MNIST peu profond, les deux convergent ;
        ReLU est légèrement plus rapide sur les architectures $\geq 3$ couches.
  \item \textbf{Tanh} : meilleur que sigmoid grâce à sa symétrie centrée en 0,
        ce qui améliore le conditionnement du gradient.
\end{itemize}

\subsection{Résultats}

\begin{table}[H]
\centering
\caption{Performances sur MNIST (100 itérations, lr=0{,}1, batch=256).}
\begin{tabular}{lrrrr}
\toprule
\textbf{Modèle} & \textbf{Params} & \textbf{Acc. Train} & \textbf{Acc. Test} & \textbf{Erreur Test} \\
\midrule
Linéaire               & 7\,850   & 93{,}06\,\% & 92{,}47\,\% & 7{,}53\,\% \\
MLP H=1 [64] sigmoid   & 50\,890  & 97{,}21\,\% & 96{,}52\,\% & 3{,}48\,\% \\
MLP H=2 [64,32] sigmoid & 52\,650 & 97{,}69\,\% & 96{,}85\,\% & 3{,}15\,\% \\
\bottomrule
\end{tabular}
\label{tab:mnist}
\end{table}

Le modèle linéaire atteint déjà $92{,}47\,\%$ — MNIST est un problème relativement simple
grâce au fond blanc uniforme et aux chiffres centrés. Le MLP H=2 gagne $+4{,}4$ points
en introduisant deux niveaux de représentation non linéaire.

\subsection{Analyse des erreurs et visualisation ACP}

Sur les $10\,000$ images de test, 238 sont mal classées simultanément par les trois modèles.
Ces erreurs se concentrent sur des paires ambiguës : $(3,8)$, $(4,9)$, $(5,6)$.

La projection ACP en 2 dimensions (figure~\ref{fig:mnist_acp}) confirme cette observation :
les classes 0 et 1 sont très bien séparées (formes distinctives),
tandis que 3/8 et 4/9 se chevauchent significativement.
Ce chevauchement est intrinsèque aux données et ne peut être réduit
qu'en exploitant la structure spatiale (ce que le MLP ne fait pas).

% ============================================================
\section{Partie 2 — CIFAR-10 : objets naturels en couleur}
% ============================================================

\subsection{Description du jeu de données}

CIFAR-10 contient $50\,000$ images $32\times32$ pixels RGB, réparties sur 10 classes
d'objets naturels (avion, cheval, chien, etc.).
L'espace de représentation est $\mathbb{R}^{3\,072}$ (couleur) ou $\mathbb{R}^{1\,024}$ (gris).

\textbf{Pourquoi CIFAR-10 est plus difficile que MNIST :}
\begin{itemize}[leftmargin=*]
  \item Fonds complexes et variés (vs fond blanc uniforme MNIST)
  \item Translations, rotations, occlusions fréquentes
  \item Couleur ajoute de la dimensionnalité sans garantir d'aide discriminante
  \item Pas d'invariance de translation dans un MLP : un avion en haut à gauche
        et un avion en bas à droite activent des poids totalement différents
\end{itemize}

\subsection{Modèles linéaires et MLP — limitations}

\begin{table}[H]
\centering
\caption{Performances des modèles denses sur CIFAR-10 (50 itérations, lr=0{,}1).}
\begin{tabular}{lccc}
\toprule
\textbf{Modèle} & \textbf{Entrée} & \textbf{Acc. Test} & \textbf{Remarque} \\
\midrule
Linéaire (gris)    & $1\,024$ & 28{,}2\,\% & Gabarit moyen par classe \\
MLP [128] (gris)   & $1\,024$ & 34{,}2\,\% & Non-linéarité insuffisante \\
Linéaire (couleur) & $3\,072$ & 40{,}0\,\% & Couleur légèrement aidante \\
MLP [128] (couleur)& $3\,072$ & 43{,}3\,\% & Plus de features, mais pas de structure \\
\bottomrule
\end{tabular}
\label{tab:cifar_dense}
\end{table}

\noindent Le MLP couleur ($43{,}3\,\%$) surpasse à peine le MLP gris ($34{,}2\,\%$) malgré
$3\times$ plus d'entrées : la couleur brute non structurée n'aide pas le MLP autant qu'on
pourrait l'espérer.

\subsection{Opération de convolution 2D}

La convolution répond au problème d'invariance spatiale en appliquant le même filtre
$K \in \mathbb{R}^{k \times k}$ en toute position de l'image :
\[
  (X * K)[i,j] = \sum_{u=0}^{k-1}\sum_{v=0}^{k-1} X[i+u,\, j+v]\cdot K[u,v]
\]
\textbf{Partage de poids :} un seul filtre $K$ (par exemple $3\times3 = 9$ paramètres)
détecte la même structure locale partout dans l'image — vs le MLP qui alloue un poids
distinct à chaque position. Cela réduit drastiquement le nombre de paramètres et confère
l'invariance de translation.

\subsection{Filtres K1–K6 du cours}

Six filtres canoniques illustrent ce que les premières couches apprennent :

\begin{table}[H]
\centering
\caption{Filtres $3\times3$ du cours et leur effet (figure~\ref{fig:filtres}).}
\begin{tabular}{llp{7cm}}
\toprule
\textbf{Filtre} & \textbf{Type} & \textbf{Effet mathématique} \\
\midrule
K1 & Flou (moyenneur)       & $K_1 = \frac{1}{9}\mathbf{1}_{3\times3}$ — supprime le bruit haute fréquence \\
K2 & Netteté (laplacien)    & $K_2[1,1]=5$, voisins $=-1$ — amplification des gradients locaux \\
K3 & Bords horizontaux      & Détecte $\partial I / \partial y$ (transitions verticales) \\
K4 & Bords verticaux        & Détecte $\partial I / \partial x$ (transitions horizontales) \\
K5 & Sobel horizontal       & Approximation robuste de $\partial I / \partial x$ \\
K6 & Diagonal               & Détecte les transitions en diagonale \\
\bottomrule
\end{tabular}
\end{table}

Les premières couches d'un CNN entraîné apprennent automatiquement des variantes
de ces filtres — c'est ce que les visualisations de Zeiler \& Fergus (2014) ont démontré.

\subsection{Architecture CNN PyTorch}

L'architecture convolutive implémentée suit le schéma du cours :

\begin{center}
\begin{tabular}{lll}
\toprule
\textbf{Couche} & \textbf{Dimensions sortie} & \textbf{Paramètres} \\
\midrule
Conv1($3\to64$, $3\times3$) + ReLU         & $(64, 32, 32)$ & $3\times3\times3\times64 = 1\,728$ \\
Conv2($64\to64$, $3\times3$) + ReLU + Pool & $(64, 16, 16)$ & $3\times3\times64\times64 = 36\,864$ \\
Conv3($64\to64$, $3\times3$) + ReLU + Pool & $(64, 8, 8)$   & $36\,864$ \\
Conv4($64\to64$, $3\times3$) + ReLU        & $(64, 8, 8)$   & $36\,864$ \\
Dropout($p=0{,}3$) + FC($4096\to10$)       & $10$           & $40\,970$ \\
\midrule
\textbf{Total}                              &               & $\mathbf{\approx 152\,778}$ \\
\bottomrule
\end{tabular}
\end{center}

\textbf{Max Pooling :} réduit la résolution par $2\times2$ en prenant le maximum local.
Effet dual : réduction de la dimension computationnelle + légère invariance de translation.

\textbf{Dropout($p=0{,}3$) :} désactive chaque neurone avec probabilité $p$ durant
l'entraînement (avec mise à l'échelle inverted dropout $1/(1-p)$).
Force des représentations redondantes, réduit l'overfitting.

\subsection{Résultats CNN et analyse de l'overfitting}

\begin{table}[H]
\centering
\caption{CNN PyTorch sur CIFAR-10 (Adam lr=0{,}001, augmentation, early stopping patience=6).}
\begin{tabular}{lccc}
\toprule
\textbf{Modèle} & \textbf{Acc. Train} & \textbf{Acc. Test} & \textbf{Gain vs MLP} \\
\midrule
MLP [128] couleur & — & 43{,}3\,\% & — \\
CNN 4 conv.       & 83{,}4\,\% & \textbf{80{,}8\,\%} & +37{,}5 pts \\
\bottomrule
\end{tabular}
\label{tab:cifar_cnn}
\end{table}

L'écart train/test de $2{,}6$ points traduit un \textbf{faible overfitting},
signe positif de l'efficacité du Dropout et de l'augmentation de données.

\textbf{Augmentation de données :} flip horizontal, crop aléatoire ($\pm4$ px), variation
de luminosité et de contraste. Ces transformations augmentent artificiellement la diversité
du dataset sans modifier les étiquettes, régularisant implicitement le modèle.

Tableau de référence scientifique pour CIFAR-10 :
\textit{Conv. Deep Belief Network} 78{,}9\,\%,
\textit{Maxout Networks} 90{,}6\,\%,
\textit{Vision Transformer (ViT)} 99{,}5\,\%.

% ============================================================
\section{Partie 3 — CBIS-DDSM : détection de cancers du sein}
% ============================================================

\subsection{Contexte médical et description du dataset}

Le dataset CBIS-DDSM (\textit{Curated Breast Imaging Subset of DDSM}) contient des
mammographies numériques annotées pour la classification bénin/malin de masses mammaires.

\textbf{Caractéristiques clés :}
\begin{itemize}[leftmargin=*]
  \item Distribution quasi-équilibrée : $48{,}3\,\%$ malins — déséquilibre faible mais traité
  \item Taille limitée : $1\,318$ images d'entraînement, $378$ de test
  \item Format : crops $224\times224$ px centrés sur la masse, niveaux de gris
  \item Enjeu clinique : un faux négatif (cancer manqué) peut être fatal
\end{itemize}

\textbf{Métriques médicales :} l'accuracy globale est insuffisante. Les métriques pertinentes sont :
\begin{align*}
  \text{Sensibilité} &= \frac{TP}{TP+FN} & \text{(cancers détectés)} \\
  \text{Spécificité} &= \frac{TN}{TN+FP} & \text{(bénins correctement classés)}
\end{align*}

\subsection{Gestion du déséquilibre de classes}

Même quasi-équilibré, nous appliquons des poids inversement proportionnels à la fréquence :
\[
  w_c = \frac{N}{2\,n_c}, \quad c \in \{0,1\}
\]
Ces poids sont intégrés dans la \texttt{CrossEntropyLoss} via
$\mathcal{L}_{CE}^w = -\sum_k w_k\,y_k\ln P_k$.
Cela pénalise plus fortement les erreurs sur la classe sous-représentée.

\subsection{Prétraitement et augmentation médicale}

Les mammographies originales ($\sim 4\,000$ px) sont réduites à $224\times224$ via
des crops centrés sur la masse annotée.
L'augmentation médicale est strictement contrôlée :

\begin{table}[H]
\centering
\caption{Augmentation médicale : transformations appliquées et justification.}
\begin{tabular}{lll}
\toprule
\textbf{Transformation} & \textbf{Paramètres} & \textbf{Justification} \\
\midrule
Flip horizontal       & $p=0{,}5$          & Masse peut apparaître des deux côtés \\
Rotation              & $\pm 15^\circ$      & Masses sans orientation fixe \\
Zoom                  & $\pm12\,\%$         & Variabilité de numérisation \\
Luminosité            & $\pm25\,\%$         & Calibrations scanner différentes \\
Contraste             & $\pm25\,\%$         & Idem \\
\textbf{Flip vertical} & \textbf{Non appliqué} & Orientation anatomique haut/bas significative \\
\bottomrule
\end{tabular}
\end{table}

Rotation et zoom sont combinés en une seule transformation affine vectorisée
(\texttt{F.affine\_grid} / \texttt{F.grid\_sample}) pour éviter deux interpolations
successives qui dégraderaient les micro-textures diagnostiques.

\subsection{Gestion mémoire GPU — accumulation de gradients}

L'entraînement sur des images $224\times224$ avec un batch de 32 nécessite
$\sim$600 MB de VRAM, dépassant la capacité disponible. La solution est l'accumulation
de gradients : on accumule les gradients sur $K$ micro-batches avant la mise à jour :
\[
  \nabla_{\text{accumulé}} = \frac{1}{K}\sum_{k=1}^{K} \nabla_{\mathcal{L}(\mathcal{B}_k)}
  \approx \nabla_{\mathcal{L}(\mathcal{B}_{\text{eff}})}, \quad
  \mathcal{B}_{\text{eff}} = K \cdot \mathcal{B}_k
\]
On utilise batch physique $= 8$ et $K = 4$, soit un batch effectif de $32$.

\textbf{Important :} BatchNorm calcule ses statistiques sur le batch \emph{physique} de 8,
pas sur le batch effectif. Ce point est discuté en détail en section~\ref{sec:ablation}.

\subsection{Modèles implémentés}

\subsubsection*{MLP profond ReLU [2048, 1024, 512, 256]}

Les images sont aplaties en vecteurs de $224\times224 = 50\,176$ valeurs.
L'architecture en entonnoir progressif vise à capturer des interactions complexes :
\[
  50\,176 \xrightarrow{ReLU} 2048 \xrightarrow{ReLU} 1024
  \xrightarrow{ReLU} 512 \xrightarrow{ReLU} 256 \xrightarrow{softmax} 2
\]
ReLU est choisi pour ses propriétés de gradient constant sur les activations positives —
indispensable sur 4 couches cachées pour éviter la disparition du gradient.

\subsubsection*{CNN from scratch (CNNBinaire)}

4 blocs Conv+BatchNorm+ReLU+MaxPool réduisent progressivement la résolution :
\[
  (1,224,224) \to (32,112,112) \to (64,56,56) \to (128,28,28) \to (128,14,14)
\]
suivis d'un classificateur dense avec Dropout(0.5 / 0.3).

L'initialisation de He (Kaiming, 2015) est appliquée explicitement à toutes les couches :
\[
  \sigma_W = \sqrt{\frac{2}{\text{fan\_out}}}
\]
Cette formule compense la réduction de variance de ReLU ($\mathbf{E}[\max(0,z)^2] = \frac{1}{2}\mathbf{E}[z^2]$
pour $z$ symétrique), garantissant une propagation stable du signal.

\subsubsection*{ResNet18 transfer learning}

ResNet18 pré-entraîné sur ImageNet ($1{,}2$M images) est adapté :
\begin{enumerate}[leftmargin=*]
  \item \texttt{conv1} recréé pour 1 canal gris (réinitialisé, ré-appris)
  \item \texttt{layer1}, \texttt{layer2} gelés — préservent les features universelles
        (bords, textures) apprises sur ImageNet, démontrées universelles par Zeiler \& Fergus (2014)
  \item \texttt{layer3}, \texttt{layer4}, \texttt{fc} fine-tunés avec lr=$5\times10^{-5}$
        pour spécialiser aux textures des masses mammaires
\end{enumerate}
Paramètres entraînables : $10\,793\,410$.

\subsection{Calibration du seuil de décision}

Le seuil $0{,}5$ par défaut n'est pas optimal en médecine.
On cherche le seuil $s^*$ maximisant l'accuracy sous contrainte de sensibilité :
\[
  s^* = \underset{s \in [0{,}05,\,0{,}95]}{\arg\max}\;\text{Acc}(s)
  \quad \text{sous contrainte : } \text{Sensibilité}(s) \geq 80\,\%
\]

\subsection{Résultats comparatifs}

\begin{table}[H]
\centering
\caption{Performances sur CBIS-DDSM ($224\times224$, seuil par défaut 0,50).}
\begin{tabular}{lcccc}
\toprule
\textbf{Modèle} & \textbf{Acc.} & \textbf{Sensib.} & \textbf{Spécif.} & \textbf{FN (cancers)} \\
\midrule
MLP [2048,1024,512,256] ReLU & 57{,}9\,\% & 64{,}6\,\% & 53{,}7\,\% & 52 / 147 \\
CNN from scratch             & 61{,}9\,\% & 39{,}5\,\% & 76{,}2\,\% & 89 / 147 \\
ResNet18 transfer            & \textbf{65{,}9\,\%} & 72{,}8\,\% & 61{,}5\,\% & 40 / 147 \\
\bottomrule
\end{tabular}
\label{tab:ddsm}
\end{table}

\begin{table}[H]
\centering
\caption{Résultats après calibration du seuil (sensibilité cible $\geq 80\,\%$).}
\begin{tabular}{lcccc}
\toprule
\textbf{Modèle} & \textbf{Seuil} & \textbf{Sensib.} & \textbf{Spécif.} & \textbf{FN réduits} \\
\midrule
MLP   & 0{,}38 & 84{,}4\,\% & 33{,}8\,\% & 23 / 147 \\
CNN   & 0{,}34 & 83{,}7\,\% & 26{,}8\,\% & 24 / 147 \\
ResNet18 & \textbf{0{,}42} & 81{,}6\,\% & 52{,}4\,\% & 27 / 147 \\
\bottomrule
\end{tabular}
\label{tab:ddsm_cal}
\end{table}

\textbf{Observation clinique :} ResNet18 avec seuil calibré atteint 81{,}6\,\% de sensibilité
avec une spécificité de 52{,}4\,\% — le meilleur équilibre.
À titre de comparaison, la littérature cite ResNet-50 à 85--88\,\% sur des datasets
bien plus larges ($\geq 5\,000$ images avec augmentation massive).

\subsection{Étude d'ablation}
\label{sec:ablation}

\subsubsection*{Protocole}

5 variantes du même CNN sont entraînées sur CBIS-DDSM ($128\times128$, 25 époques),
chacune activant une technique supplémentaire. L'initialisation de Kaiming et l'augmentation
médicale sont communes à toutes les configurations.

\subsubsection*{Résultats et analyse}

\begin{table}[H]
\centering
\caption{Étude d'ablation (CBIS-DDSM $128\times128$, batch physique=8, effectif=32).}
\begin{tabular}{lccccrr}
\toprule
\textbf{Config.} & \textbf{BN} & \textbf{Drop.} & \textbf{WD} & \textbf{Perte} &
\textbf{Acc.} & \textbf{$\mathcal{L}_{\text{init}}$} \\
\midrule
1. Baseline (aucune régul.)   & $\times$ & $\times$ & 0       & CE  & \textbf{63{,}8\,\%} & 0{,}95 \\
2. + BatchNorm                & \checkmark & $\times$ & 0     & CE  & 57{,}7\,\% & 31{,}6 \\
3. + Dropout seul             & $\times$ & \checkmark & 0     & CE  & 51{,}6\,\% & 1{,}79 \\
4. + BN + Dropout + WD        & \checkmark & \checkmark & $10^{-4}$ & CE & 58{,}7\,\% & 70{,}2 \\
5. Complet avec BCE           & \checkmark & \checkmark & $10^{-4}$ & BCE & 47{,}9\,\% & 63{,}9 \\
\bottomrule
\end{tabular}
\label{tab:ablation}
\end{table}

\paragraph{Analyse.}

\textbf{(1) Le baseline gagne (63{,}8\,\%)} — avec $1\,318$ images et 25 époques,
le réseau n'a pas le temps de sur-apprendre. La régularisation n'a pas d'overfitting
à corriger ; elle ralentit simplement la convergence.

\textbf{(2) BatchNorm instable (perte initiale $31{,}6$)} — BN normalise par statistiques
de mini-batch :
\[
  \hat{x}_i = \frac{x_i - \mu_\mathcal{B}}{\sqrt{\sigma_\mathcal{B}^2 + \varepsilon}}
\]
Avec seulement 8 images, $\mu_\mathcal{B}$ et $\sigma_\mathcal{B}^2$ sont des estimations
très bruitées. La perte démarre à $31{,}6$ au lieu de $\sim0{,}7$ pour un réseau correctement
conditionné. \textit{Note :} l'accumulation de gradients (batch effectif $= 32$) ne résout
\textbf{pas} ce problème — BN calcule ses statistiques sur le batch physique.

\textbf{(3) Dropout seul cause du sous-apprentissage} — sans BN pour stabiliser
les activations, Dropout(0.5) coupe 50\,\% des connexions à chaque forward pass.
Le gradient moyen reçu par chaque poids est divisé par deux,
insuffisant avec 1\,318 images et batch de 8. Early stopping à l'époque 8.

\textbf{(4) Bug identifié — weight decay sur les paramètres BN} — dans les configs 4 et 5,
\texttt{Adam(model.parameters(), weight\_decay=$10^{-4}$)} appliquait la pénalité L2
aux paramètres $\gamma$ et $\beta$ de BatchNorm. Or $\gamma$ régularisé par L2 tend vers 0,
supprimant progressivement la normalisation. Ce bug est désormais corrigé en utilisant
deux groupes de paramètres (BN exclus du weight decay).

\textbf{(5) BCE sous le hasard (47{,}9\,\%)} — l'instabilité BN + le \texttt{pos\_weight}
pousse le modèle vers une solution dégénérée : prédire ``Malin'' par défaut
(sensibilité 66\,\% mais spécificité 36\,\%).

\subsubsection*{Perspectives d'amélioration}

Deux techniques adressent directement les causes d'instabilité identifiées dans l'ablation.

\textbf{1 — BN à momentum faible ($m = 0{,}01$) :}
\[
  \mu_{\text{run}} \leftarrow (1-m)\,\mu_{\text{run}} + m\,\mu_\mathcal{B}
\]
Avec $m = 0{,}01$ (au lieu de $m = 0{,}1$ par défaut), les running stats convergent
lentement : 100 batches pour converger à 63\,\% au lieu de 10 batches.
Cela rend BN robuste aux petits batches en ne surpondérant pas les statistiques
instantanées bruitées.

\textbf{2 — MixUp ($\alpha = 0{,}4$) :}
\[
  \tilde{x} = \lambda\,x_i + (1-\lambda)\,x_j, \quad
  \tilde{y} = \lambda\,y_i + (1-\lambda)\,y_j, \quad
  \lambda \sim \mathrm{Beta}(0{,}4,\,0{,}4)
\]
Force la surface de décision à être linéaire \emph{entre} les exemples d'entraînement.
Régularise via les données (pas via l'architecture), sans perturber les gradients.
La perte MixUp est :
$\mathcal{L} = \lambda\,\mathcal{L}(f(\tilde{x}), y_i) + (1-\lambda)\,\mathcal{L}(f(\tilde{x}), y_j)$

\begin{table}[H]
\centering
\caption{Perspectives — configurations 6 et 7 (CBIS-DDSM $128\times128$, 25 époques).
         Ces deux configurations corrigent l'instabilité BN identifiée en section~\ref{sec:ablation}.}
\small
\begin{tabular}{lccp{4.5cm}c}
\toprule
\textbf{Config.} & \textbf{BN mom.} & \textbf{MixUp} & \textbf{Objectif} & \textbf{Acc.} \\
\midrule
4. Référence (BN+Drop+WD)              & 0{,}10 & Non & Baseline régularisé & 59{,}5\,\% \\
6. + BN momentum $= 0{,}01$            & 0{,}01 & Non & Stabiliser stats.\ BN small batch & 60{,}3\,\% \\
7. + BN $m{=}0{,}01$ + MixUp $\alpha{=}0{,}4$ & 0{,}01 & Oui & Régularisation combinée & 51{,}6\,\% \\
\bottomrule
\end{tabular}
\label{tab:perspectives}
\end{table}

\noindent\textbf{Résultats obtenus.}
La config~6 atteint $60{,}3\,\%$ (vs $59{,}5\,\%$ pour la config~4), gain marginal :
le momentum réduit ne résout pas l'instabilité fondamentale — la perte initiale
reste $\times 75$ (config~6~: $75{,}3$ vs config~4~: $73{,}2$).
La config~7 ($51{,}6\,\%$) est la moins performante en accuracy, mais c'est la deuxième
configuration la plus équilibrée (sensibilité~$40{,}1\,\%$, spécificité~$58{,}9\,\%$)
après la baseline~1 — MixUp régularise effectivement la surface de décision.
La conclusion est que le problème de fond reste le batch physique de~8 : aucune
technique appliquée à BN seul ne suffit tant que cette contrainte n'est pas levée.

% ============================================================
\section{Conclusion et synthèse}
% ============================================================

Ce projet illustre comment la complexité des données impose des choix architecturaux
et mathématiques de plus en plus sophistiqués.

\begin{table}[H]
\centering
\caption{Synthèse des trois parties.}
\small
\begin{tabular}{lp{3.2cm}p{3.2cm}p{3cm}c}
\toprule
\textbf{Dataset} & \textbf{Contrainte principale} & \textbf{Solution clé} & \textbf{Meilleur modèle} & \textbf{Acc.} \\
\midrule
MNIST     & Classification simple & MLP H=2 sigmoid & MLP [64,32] sigmoid & 96{,}9\,\% \\
CIFAR-10  & Invariance spatiale   & Convolutions    & CNN 4 blocs         & 80{,}8\,\% \\
CBIS-DDSM & Petit dataset, FN critiques & Transfer + calibration & ResNet18 & 65{,}9\,\% \\
\bottomrule
\end{tabular}
\end{table}

\textbf{Leçons mathématiques transversales :}
\begin{enumerate}[leftmargin=*]
  \item \textbf{L'architecture doit correspondre à la structure des données.}
        Les convolutions exploitent la localité spatiale — un MLP sur CIFAR-10 perd
        toute information de position relative des pixels.
  \item \textbf{La régularisation n'est bénéfique qu'en présence d'overfitting.}
        Sur $1\,318$ images en 25 époques, le Dropout et le BN nuisent à la convergence
        car il n'y a pas encore de sur-apprentissage à corriger.
  \item \textbf{BatchNorm requiert des batches physiques $\geq 32$.}
        L'accumulation de gradients simule un grand batch pour les gradients mais pas
        pour les statistiques de normalisation — un piège fréquent en pratique.
  \item \textbf{Le seuil de décision est un paramètre du modèle, pas une constante.}
        En médecine, calibrer le seuil permet de passer de 40\,\% à 84\,\% de sensibilité
        au prix d'une baisse de spécificité — un compromis explicite et maîtrisé.
  \item \textbf{Le transfer learning compense le manque de données.}
        ResNet18 pré-entraîné sur ImageNet apporte des features universelles
        (bords, textures) qui restent pertinentes même en imagerie médicale.
\end{enumerate}

% ============================================================
\clearpage
\section*{Annexes — Figures}
% ============================================================
\addcontentsline{toc}{section}{Annexes}

\begin{figure}[H]
    \centering
    \includegraphics[width=0.55\textwidth]{sorties/mnist/convergence.png}
    \caption{Convergence de la log-loss sur l'ensemble d'entraînement MNIST (3 modèles).}
    \label{fig:mnist_conv}
\end{figure}

\begin{figure}[H]
    \centering
    \includegraphics[width=0.55\textwidth]{sorties/mnist/acp.png}
    \caption{Projection ACP 2D des données MNIST.
             Les classes 0 et 1 sont bien séparées ;
             les paires (3,8), (4,9) se chevauchent, expliquant les erreurs résiduelles.}
    \label{fig:mnist_acp}
\end{figure}

\begin{figure}[H]
    \centering
    \includegraphics[width=0.6\textwidth]{sorties/mnist/erreurs_communes.png}
    \caption{Exemples de chiffres ambigus mal classés simultanément par les trois modèles.}
    \label{fig:mnist_erreurs}
\end{figure}

\begin{figure}[H]
    \centering
    \includegraphics[width=0.95\textwidth]{sorties/cifar/filtres_convolution.png}
    \caption{Effet des 6 filtres K1--K6 sur une image en niveaux de gris.
             De gauche à droite : original, flou, netteté, bords horizontaux,
             bords verticaux, Sobel, diagonal.}
    \label{fig:filtres}
\end{figure}

\begin{figure}[H]
    \centering
    \includegraphics[width=0.6\textwidth]{sorties/cifar/convergence.png}
    \caption{Convergence train/validation du CNN PyTorch sur CIFAR-10.}
    \label{fig:cifar_conv}
\end{figure}

\begin{figure}[H]
    \centering
    \includegraphics[width=0.6\textwidth]{sorties/cifar/acp.png}
    \caption{Projection ACP 2D de CIFAR-10. Le chevauchement important entre toutes les classes
             illustre pourquoi les modèles linéaires et MLP peinent sur ce dataset.}
    \label{fig:cifar_acp}
\end{figure}

\begin{figure}[H]
    \centering
    \begin{subfigure}[b]{0.48\textwidth}
        \includegraphics[width=\textwidth]{sorties/ddsm/mlp_confusion_05.png}
        \caption{MLP — seuil 0,5 (sensi. 64,6\,\%)}
    \end{subfigure}
    \hfill
    \begin{subfigure}[b]{0.48\textwidth}
        \includegraphics[width=\textwidth]{sorties/ddsm/mlp_confusion_cal.png}
        \caption{MLP — seuil calibré 0,38 (sensi. 84,4\,\%)}
    \end{subfigure}
    \caption{Matrices de confusion du MLP sur CBIS-DDSM.}
    \label{fig:ddsm_mlp}
\end{figure}

\begin{figure}[H]
    \centering
    \begin{subfigure}[b]{0.48\textwidth}
        \includegraphics[width=\textwidth]{sorties/ddsm/cnn_confusion_05.png}
        \caption{CNN from scratch — seuil 0,5}
    \end{subfigure}
    \hfill
    \begin{subfigure}[b]{0.48\textwidth}
        \includegraphics[width=\textwidth]{sorties/ddsm/resnet_confusion_cal.png}
        \caption{ResNet18 — seuil calibré 0,42}
    \end{subfigure}
    \caption{Matrices de confusion CNN et ResNet18 sur CBIS-DDSM.}
    \label{fig:ddsm_cnn}
\end{figure}

\begin{figure}[H]
    \centering
    \begin{subfigure}[b]{0.48\textwidth}
        \includegraphics[width=\textwidth]{sorties/ddsm/ablation_train_convergence.png}
        \caption{Perte entraînement}
    \end{subfigure}
    \hfill
    \begin{subfigure}[b]{0.48\textwidth}
        \includegraphics[width=\textwidth]{sorties/ddsm/ablation_val_convergence.png}
        \caption{Perte validation}
    \end{subfigure}
    \caption{Ablation study — convergence des 7 configurations sur CBIS-DDSM ($128\times128$).
             La config~2 (BN seul) montre l'instabilité liée au batch physique de 8.
             Les configs~6 et~7 (perspectives) devraient montrer une perte initiale plus stable.}
    \label{fig:ablation}
\end{figure}

\end{document}
